% Abhakseite „Das kann ich“ – Zone nur im Gesamt (\mitzone)
%% TODO Ich-kann-Sätze: Platzhalter wie in den Aufgabendateien
\begin{abhakseite}
\ifdefined\mitzone
\abhakgruppe{Kennst du schon}
\abhak{1}{Punkte im Koordinatensystem eintragen und ablesen (erster Quadrant)}
\abhak{2}{Multiplizieren und Dividieren mit Dezimalzahlen (Preis geteilt durch Stückzahl, Preis mal Stückzahl)}
\abhak{3}{Vielfache und Teiler erkennen (zwölf ist das Dreifache von vier)}
\abhak{4}{Bruchteil einer Größe (die Hälfte, ein Viertel von 12 €)}
\abhak{5}{Einheiten umrechnen (Cent ↔ Euro, Minuten ↔ Stunden, m ↔ km)}
\abhak{6}{Vielfache und Teiler erkennen (zwölf ist das Dreifache von vier) – Fehler finden}
\abhak{7}{Vielfache und Teiler erkennen (zwölf ist das Dreifache von vier) – selbst rechnen}
\fi
\abhakgruppe{Einheit 1 · Zuordnungen darstellen}
\abhak{8}{Darstellen}
\abhak{9}{Achseneinteilung}
\abhak{10}{Formel aus Tabelle (y = 4 · x)}
\abhak{11}{Darstellen – Fehler finden, Begründen, Darstellungswechsel, Anwendung}
\abhakgruppe{Einheit 2 · Proportionale Zuordnungen und Dreisatz}
\abhak{12}{Was ist hier eine Portion?}
\abhak{13}{Minitabelle anlegen}
\abhak{14}{Pfeile beidseitig}
\abhak{15}{Dreisatz proportional}
\abhak{16}{fester Faktor angeben (Preis je Einheit), Gleichung y = k · x}
\abhak{17}{Graph zeichnen (Ursprungsgerade) und Werte ablesen}
\abhak{18}{Verhältnisgleichung aufstellen (Vorrat)}
\abhak{19}{Dreisatz proportional – Fehler finden, Begründen, Darstellungswechsel, Anwendung}
\abhakgruppe{Einheit 3 · Antiproportionale Zuordnungen}
\abhak{20}{Antiproportional}
\abhak{21}{Graph als fallende Kurve, Werte ablesen}
\abhak{22}{Antiproportional – Fehler finden, Begründen, Darstellungswechsel, Anwendung}
\abhakgruppe{Einheit 4 · Zuordnungstypen erkennen und anwenden}
\abhak{23}{Nullwert prüfen}
\abhak{24}{Erkennen}
\abhak{25}{Rate}
\abhak{26}{Gegenbeispiele (Alter und Größe, Tarif mit Grundgebühr)}
\abhak{27}{Erkennen – Fehler finden, Begründen, Darstellungswechsel, Anwendung}
\end{abhakseite}
